\subsection{Tamaño de Métodos y Archivos}

\subsubsection{Regla: máximo de 30 líneas por método y 300 por archivo}

\textbf{Regla:} El cuerpo de un método no supera 30 líneas, sin contar las llaves que lo abren y lo cierran. Un archivo de código no supera 300 líneas. Al superar cualquiera de los dos límites, se extrae la parte con una responsabilidad propia a un método privado o a una clase nueva (McConnell, 2004; Martin, 2008). El código generado por herramientas, como \texttt{Strings.Designer.cs}, queda fuera de este límite.

\textbf{Descripción:} McConnell señala que las rutinas largas acumulan responsabilidades y concentran más errores por línea a partir de cierto tamaño (McConnell, 2004), y Martin recomienda que una función haga una sola cosa (Martin, 2008). Un límite explícito convierte esas recomendaciones en un criterio que se puede revisar en cada solicitud de integración.

\textbf{Código correcto:}
\begin{lstlisting}[style=csharpstyle]
public void StartMatch(MatchSetup setup)
{
    PlacePlayers(setup.Players);
    ShuffleDeck(setup.Deck);
    AnnounceFirstTurn(setup.Players);
}
\end{lstlisting}

\textbf{Código incorrecto:} un método \texttt{StartMatch} de 60 líneas que coloca a los jugadores, baraja el mazo y anuncia el primer turno dentro del mismo cuerpo.
